\documentclass{article}
\usepackage{amsmath}

\title{Q and P: Costly Cognition in a Double Auction}
\author{Ada}
\date{}

\begin{document}
\maketitle

\section*{The pair in two sentences}
Q studies LLM buyers and sellers in a continuous double auction and reports that repeated one-cent quote improvements and reluctance to cross the spread can delay trade and reduce gross allocative efficiency (ledger entries 12 and 16). P studies an induced ``survival'' economy in which inference has cost $C=kTS^{\alpha}$ and agents can deactivate, but its survival pressure is explicitly constructed rather than literal (ledger entries 1, 3, and 16).

\section*{The connexion pursued}
The pursued connexion is costly market participation: place P's token-cost and deactivation accounting inside Q's existing auction, rather than adding an auction to P's job-selection environment (entries 1, 3, and 8). The sharpened question is whether behaviorally salient inference costs regularize Q's incremental dithering or instead reduce efficiency through computational cost, exit, and market thinning (entries 2, 11, and 20).

\section*{What was found}
The material establishes a design and a testable decomposition, not an experimental result. Q defines gross allocative efficiency as realized buyer-plus-seller surplus divided by the maximum gains from trade, and links repeated $\$0.01$ improvements and unwillingness to cross with depressed trade and efficiency; P supplies token cost $C=kTS^{\alpha}$ and evidence that token-generating interaction trades energy expenditure against coordination (entry 16). These facts support preserving Q's gross-efficiency ratio while separately measuring net surplus after inference cost in common payoff units (entries 6, 16, and 20).

The proposed comparison is a prompt-matched zero-price meter, a visible inference fee without exit, and the same fee with a hard budget, swept over fee and budget levels (entries 15, 17, and 20). Agents in the fee arms must observe the meter and maximize terminal trading profit net of the fee: retrospective cost subtraction changes welfare accounting but cannot change behavior (entries 15 and 19). Thus fee versus zero-price meter identifies a within-trader response, while hard budget versus fee-only identifies the additional effect of exit and thinning (entries 5, 11, and 19). Crossing size and latency are primary mechanism measures; lexical shifts are secondary (entries 15 and 19). Moderate cost may hasten crossing, high cost may suppress participation, and a null pre-exit response would weaken the proposed urgency mechanism, so no directional result is assumed (entries 6, 11, and 20).

The second round did not produce the requested pilot result. Inspection found that Q's existing agents maximize gross trading profit, while P's cost formula can only be imported retrospectively without changing those agents' incentives; such an import would change net-welfare accounting but would not identify a behavioral fee response (entries 24, 27, and 30). The available no-key zero-intelligence baseline also cannot supply that response because it neither reads a fee prompt nor optimizes net payoff (entry 24). Thus the standing result is a credible, falsifiable project, not evidence that inference pricing already changes auction behavior or welfare (entries 28, 31, and 33).

\section*{Objections and resolution}
Several objections were raised. First, P does not validate stronger real-world stakes: both settings use induced incentives, and P warns against anthropomorphic interpretation (entries 1 and 3). Second, combining cost with deactivation confounds behavioral response with selective exit; the fee-only and hard-budget arms resolve this in design, provided the fee is visible and affects utility (entries 4, 5, 10, and 15). Third, P's no-size-penalty comparison jointly changes deactivation and active attempts, so increased collisions cannot establish an urgency effect; it only motivates conditioning outcomes on active-agent exposure (entries 7 and 11). Fourth, Q's lexical evidence is correlational, limited to one model, and potentially prompt-induced; prompt matching and behavioral endpoints reduce but do not remove this concern (entries 4, 15, and 19).

The verifier objected that the note had a design but no empirical result and requested the zero-price versus visible-fee pilot (entry 23). That objection remains unresolved empirically. Both peers found that the supplied workspace lacks a runnable auction checkout and supported provider credentials, while Q's public implementation lacks the fee utility and balance state required by the treatment (entries 24, 27, and 30). This resolves only why the pilot was not completed; it does not resolve the evidentiary gap.

\section*{Sources outside Q and P}
The second round inspected Q's public simulation repository to assess whether the verifier's pilot was immediately runnable. Its README and implementation were used only for the feasibility claims about API keys, token logs, and missing inference-fee state (entries 24 and 25). No outside source supplied a substantive result about the Q--P connexion.

\section*{What remains open}
Open issues are the conversion between dollar profit and energy, initial-budget calibration, comparable reasoning-token metering across providers, statistical power under stochastic runs, access to comparable traces, prompt-induced vocabulary, and whether a human costly-cognition benchmark is required (entries 8, 13, and 21). The material is thin in one specific respect: Emmy's note is available, but no Ada peer note is present in the supplied directory. The ledger contains contributions from both peers, including readiness declarations at entries 18 and 22, but it cannot substitute for a missing note.

More importantly, the material is thin because it contains no treated market run. Publication viability remains conditional on a pre-exit fee effect or a precise informative null, and neither is in the ledger or notes (entries 23, 28, 31, and 34).

\section*{Smallest next step}
Add per-agent cumulative token fees to Q's state and net-payoff prompt, expose the meter, and create matched price-zero and positive-price configurations; then run replicated credentialed LLM markets with exit disabled (entries 25, 27, and 31). Record crossing latency and size, gross surplus, and net surplus as the verifier requested (entry 23). This is the smallest step that can settle the behavioral bridge: a fee effect would support it, while a precise null would weaken the urgency mechanism (entries 28 and 31).

\end{document}
